\documentclass[11pt]{article}
\usepackage[margin=1in]{geometry}
\usepackage{hyperref}
\usepackage{enumitem}
\title{Plant Care Tracker: AI Tooling Disclosure}
\author{Josh Dunlap, Aaron Fuentes, Nathan Reeves \\ CISC 450 (Spring 2026)}
\date{}
\begin{document}
\maketitle
\section*{AI Tooling Disclosure}
\subsection*{Tool used}
The development of this project's source code was assisted by \textbf{Claude
Code}.
\subsection*{Division of labor}
We owned the parts of the project that involve thinking: the
database design, the system requirements, every structural decision about
how the data is modeled, and the visual/UX direction. Claude was used
mainly to write the \emph{code} that implements those decisions.\\

\noindent In particular, \textbf{the database side of the project was primarily our work}. We designed the ER model in Lucid, decided what columns and relationships should exist, chose which deletes should cascade vs
restrict, and decided when to add new columns (e.g.\ \texttt{birthday} alongside
\texttt{approx\_age}, \texttt{image\_url} for photo uploads) as new features came up.
Claude's role on the database side was to review the schema and suggest
fixes when issues came up during the build, and to write the migration
SQL we decided on.
\subsection*{What AI helped with}
At a high level, Claude was used for:
\begin{itemize}[leftmargin=*]
  \item \textbf{Scaffolding the application:} initial Next.js project setup,
    configuration for the ORM and the deployment, env-var wiring, and CLI
    provisioning steps that are tedious but mechanical.
  \item \textbf{Writing UI components} from team-provided layouts and feature
    descriptions. We described what each page should do; Claude
    produced the React/Tailwind components, which we then reviewed.
  \item \textbf{Writing query and Server Action code} to back our specified
    requirements. Claude translated ``we need a page that shows overdue
    plants'' into actual SQL/Drizzle code; we reviewed for
    correctness.
  \item \textbf{Iterating on features} we requested during the build: the
    retro visual overhaul, separating Location and Room as independent
    attributes, the inline-create flow on the Add Plant form, the photo
    upload, and so on. Each was driven by team direction.
  \item \textbf{Debugging production and git issues,} most notably a 504-timeout problem
    caused by schema-lock contention on cold function starts. Claude
    analyzed the logs, proposed a hypothesis, and we confirmed and
    approved the fix.
  \item \textbf{\LaTeX{} scaffolds of the deliverable docs} which we completed.
\end{itemize}
\subsection*{Where human intervention was required}
\begin{itemize}[leftmargin=*]
  \item \textbf{Database design end-to-end} was our domain. Every column,
    relationship, and constraint in \texttt{schema.pdf} reflects our design,
    not Claude's.
  \item \textbf{Requirements interpretation:} the milestone-1 requirements PDF was
    given to Claude as a fixed spec; bonus features (calendar view, search,
    photo upload, retro UI) were our requested additions, not
    Claude-initiated.
  \item \textbf{Account-bound operations} Supabase project management, Vercel
    deploys, and Lucid diagram editing were all done by us.
  \item \textbf{The schema diagram} was produced in Lucid by us. Claude only
    suggested specific edits for us to apply.
  \item \textbf{The reflection document} was written by us using a \LaTeX{} template that Claude generated; the writing was done without AI assistance.
\end{itemize}
\subsection*{Statement}
All AI-generated code was reviewed by at least one team member before
being committed. The team is collectively responsible for the project's
design, correctness, and academic integrity.
\end{document}